\documentclass[12pt,a4paper]{article}

% ─── Packages ─────────────────────────────────────────────────────────────────
\usepackage[utf8]{inputenc}
\usepackage[T1]{fontenc}
\usepackage[english]{babel}
\usepackage{geometry}
\usepackage{graphicx}
\usepackage{hyperref}
\usepackage{listings}
\usepackage{xcolor}
\usepackage{booktabs}
\usepackage{longtable}
\usepackage{array}
\usepackage{tabularx}
\newcolumntype{D}{>{\raggedright\arraybackslash}p{8.5cm}}
\newcolumntype{E}{>{\raggedright\arraybackslash}p{7cm}}
\newcolumntype{F}{>{\raggedright\arraybackslash}p{6cm}}
\usepackage{fancyhdr}
\usepackage{titlesec}
\usepackage{enumitem}
\usepackage{tcolorbox}
\usepackage{amsmath}
\usepackage{lmodern}
\usepackage{parskip}
\usepackage[titles]{tocloft}

% ─── Page Geometry ────────────────────────────────────────────────────────────
\geometry{
  a4paper,
  top=2.5cm,
  bottom=2.5cm,
  left=2.5cm,
  right=2.5cm,
  headheight=15pt
}

% ─── Colors ───────────────────────────────────────────────────────────────────
\definecolor{primary}{RGB}{37, 99, 235}
\definecolor{secondary}{RGB}{15, 23, 42}
\definecolor{accent}{RGB}{99, 102, 241}
\definecolor{success}{RGB}{22, 163, 74}
\definecolor{warning}{RGB}{234, 88, 12}
\definecolor{danger}{RGB}{220, 38, 38}
\definecolor{muted}{RGB}{100, 116, 139}
\definecolor{codebg}{RGB}{248, 250, 252}
\definecolor{codeborder}{RGB}{226, 232, 240}
\definecolor{getcolor}{RGB}{59, 130, 246}
\definecolor{postcolor}{RGB}{16, 185, 129}
\definecolor{putcolor}{RGB}{245, 158, 11}
\definecolor{deletecolor}{RGB}{239, 68, 68}

% ─── Hyperref Setup ───────────────────────────────────────────────────────────
\hypersetup{
  colorlinks=true,
  linkcolor=primary,
  urlcolor=primary,
  citecolor=primary,
  pdftitle={FTP Drive API Documentation},
  pdfauthor={School/Team},
  pdfsubject={API Reference},
}

% ─── Header / Footer ──────────────────────────────────────────────────────────
\pagestyle{fancy}
\fancyhf{}
\fancyhead[L]{\small\textcolor{muted}{FTP Drive API}}
\fancyhead[R]{\small\textcolor{muted}{\leftmark}}
\fancyfoot[C]{\small\textcolor{muted}{\thepage}}
\renewcommand{\headrulewidth}{0.4pt}
\renewcommand{\footrulewidth}{0.4pt}

% ─── Section Title Styles ─────────────────────────────────────────────────────
\titleformat{\section}
  {\color{secondary}\Large\bfseries}
  {\thesection}{1em}{}[\color{primary}\titlerule]

\titleformat{\subsection}
  {\color{secondary}\large\bfseries}
  {\thesubsection}{1em}{}

\titleformat{\subsubsection}
  {\color{secondary}\normalsize\bfseries}
  {\thesubsubsection}{1em}{}

% ─── Code Listing Style ───────────────────────────────────────────────────────
\lstdefinestyle{jsonStyle}{
  backgroundcolor=\color{codebg},
  basicstyle=\ttfamily\small,
  breakatwhitespace=false,
  breaklines=true,
  captionpos=b,
  commentstyle=\color{muted},
  frame=single,
  rulecolor=\color{codeborder},
  keepspaces=true,
  keywordstyle=\color{primary}\bfseries,
  stringstyle=\color{success},
  numberstyle=\tiny\color{muted},
  numbers=none,
  showspaces=false,
  showstringspaces=false,
  tabsize=2,
  xleftmargin=5pt,
  xrightmargin=5pt,
  aboveskip=6pt,
  belowskip=6pt,
}

\lstset{style=jsonStyle}

% ─── Custom Boxes ─────────────────────────────────────────────────────────────
\tcbuselibrary{skins,breakable}

\newtcolorbox{notebox}[1][]{
  enhanced,
  colback=blue!5!white,
  colframe=primary,
  fonttitle=\bfseries,
  title=Note,
  #1
}

\newtcolorbox{warningbox}[1][]{
  enhanced,
  colback=orange!5!white,
  colframe=warning,
  fonttitle=\bfseries,
  title=Warning,
  #1
}

% ─── HTTP Method Commands ─────────────────────────────────────────────────────
\newcommand{\httpget}[1]{%
  \colorbox{getcolor!20}{\textcolor{getcolor}{\textbf{\small GET}}}~\texttt{#1}%
}
\newcommand{\httppost}[1]{%
  \colorbox{postcolor!20}{\textcolor{postcolor}{\textbf{\small POST}}}~\texttt{#1}%
}
\newcommand{\httpput}[1]{%
  \colorbox{putcolor!20}{\textcolor{putcolor}{\textbf{\small PUT}}}~\texttt{#1}%
}
\newcommand{\httpdelete}[1]{%
  \colorbox{deletecolor!20}{\textcolor{deletecolor}{\textbf{\small DELETE}}}~\texttt{#1}%
}

% ─── Misc ─────────────────────────────────────────────────────────────────────
\newcommand{\required}{\textcolor{danger}{\textbf{*}}}
\newcommand{\optional}{\textcolor{muted}{\textit{opt}}}

% ══════════════════════════════════════════════════════════════════════════════
%  DOCUMENT
% ══════════════════════════════════════════════════════════════════════════════
\begin{document}

% ─── Title Page ───────────────────────────────────────────────────────────────
\begin{titlepage}
  \centering
  \vspace*{2cm}

  {\color{primary}\rule{\linewidth}{2pt}}
  \vspace{0.5cm}

  {\Huge\bfseries\color{secondary} FTP Drive API}\\[0.4cm]
  {\Large\color{muted} Technical Reference Documentation}\\[0.2cm]
  {\normalsize\color{muted} Version 1.0 \textemdash{} June 2026}

  \vspace{0.5cm}
  {\color{primary}\rule{\linewidth}{2pt}}

  \vspace{2cm}

  \begin{tcolorbox}[colback=primary!5,colframe=primary,width=0.85\linewidth]
    \centering
    \textbf{Project:} FTP-Based Cloud Drive Storage API\\
    \textbf{Framework:} Hono (TypeScript / Node.js)\\
    \textbf{Database:} PostgreSQL via Prisma ORM\\
    \textbf{Auth:} Google OAuth 2.0 + JWT (HS256)\\
    \textbf{Storage:} Multi-site FTP (basic-ftp)\\
    \textbf{Jobs:} Inngest Background Functions
  \end{tcolorbox}

  \vfill

  {\small\color{muted}
    Confidential \textemdash{} Internal Use Only\\
    \textit{user@email.sch.id}
  }
\end{titlepage}

% ─── Table of Contents ────────────────────────────────────────────────────────
\tableofcontents
\newpage

% ══════════════════════════════════════════════════════════════════════════════
\section{Overview}
% ══════════════════════════════════════════════════════════════════════════════

\subsection{Introduction}

The \textbf{FTP Drive API} is a backend service that provides cloud-storage functionality backed by one or more FTP servers. It exposes a RESTful HTTP API under the \texttt{/v1} prefix, enabling authenticated clients to manage folders, upload/download files, share resources, and administer the system. Authentication is handled via Google OAuth 2.0 combined with stateful JWT sessions stored in the database.

\subsection{Technology Stack}

\begin{table}[h!]
\centering
\begin{tabularx}{\linewidth}{lX}
\toprule
\textbf{Component} & \textbf{Technology} \\
\midrule
Runtime      & Node.js (ESM) \\
Language     & TypeScript 5.x \\
Web Framework & Hono 4.x \\
ORM          & Prisma 7.x \\
Database     & PostgreSQL (schema: \texttt{drive}) \\
FTP Client   & basic-ftp 6.x \\
Auth         & Google OAuth 2.0 + JWT (HS256) \\
Background Jobs & Inngest 4.x \\
Validation   & class-validator + class-transformer \\
Logging      & Winston \\
Observability & OpenTelemetry (Traceloop instrumentation) \\
\bottomrule
\end{tabularx}
\caption{Technology stack}
\end{table}

\subsection{Base URL}

All endpoints listed in this document are relative to the API base URL and the version prefix \texttt{/v1}:

\begin{lstlisting}
https://<API_URL>/v1/<endpoint>
\end{lstlisting}

\subsection{Business Rules}

The following rules govern the system behaviour:

\begin{enumerate}[label=\arabic*.]
  \item There are \textbf{3 FTP service sites}, each identified by a numeric \texttt{siteId}. Files are stored on a specific site determined at upload time.
  \item Every file upload stores the actual bytes on the designated FTP site.
  \item \textbf{Delete is a hard-delete}: removing a file or folder deletes both the database record and the FTP artefact.
  \item Users can select multiple folders at once during upload, and the system will recursively create all necessary parent directories on the FTP server before writing any file.
  \item Login is exclusively via Google (\texttt{gmail}) — no local passwords are stored.
  \item Communication with the FTP server is \textbf{encrypted}.
  \item Uploaded files/folders are assigned UNIX permissions (\texttt{chmod}) matching the currently authenticated user.
  \item \textbf{Folder sharing} propagates \emph{bottom-up}: sharing a folder automatically shares all its ancestors.
  \item \textbf{File sharing} propagates \emph{top-down}: sharing a file implicitly shares its containing folder path.
\end{enumerate}

\newpage

% ══════════════════════════════════════════════════════════════════════════════
\section{Authentication}
% ══════════════════════════════════════════════════════════════════════════════

\subsection{Authentication Flow}

The API uses Google OAuth 2.0 for identity and a stateful JWT mechanism for API access control. The flow is:

\begin{enumerate}[label=\arabic*.]
  \item Client redirects the user to \texttt{GET /v1/oauth/google}.
  \item The API generates a Google OAuth consent URL and returns a redirect.
  \item Google redirects back to \texttt{GET /v1/oauth/google/callback} with an authorisation code.
  \item The API exchanges the code for tokens, fetches the Google user profile, upserts the account in the database, and issues a JWT \texttt{accessToken} and \texttt{refreshToken}.
  \item The client includes the \texttt{accessToken} in subsequent requests as a \texttt{Bearer} token.
  \item The active session is stored in the \texttt{Session} table. Logout or token expiry invalidates the session.
\end{enumerate}

\subsection{Authorization Header}

All protected endpoints require:

\begin{lstlisting}
Authorization: Bearer <accessToken>
\end{lstlisting}

The \texttt{Guard.validate()} middleware:
\begin{itemize}
  \item Verifies the JWT signature using \texttt{HS256} and the \texttt{JWT\_SECRET} env variable.
  \item Confirms the account is \texttt{ACTIVE} in the database.
  \item Confirms the session record exists and is \texttt{ACTIVE} (replay-attack protection).
  \item Injects the \texttt{account} context object into the Hono request context.
\end{itemize}

\subsection{Auth Endpoints}

\subsubsection{\httpget{/v1/oauth/google}}

Initiates Google OAuth handshake. Redirects the browser to Google's consent page.

\textbf{Auth required:} No

\textbf{Response:} \texttt{302 Redirect} $\rightarrow$ Google consent URL

\vspace{0.5em}
\subsubsection{\httpget{/v1/oauth/google/callback}}

Handles the OAuth callback from Google.

\textbf{Auth required:} No

\textbf{Query Parameters:}

\begin{tabularx}{\linewidth}{llX}
\toprule
\textbf{Parameter} & \textbf{Type} & \textbf{Description} \\
\midrule
\texttt{code} & string & OAuth authorisation code from Google \\
\texttt{state} & string & CSRF state token \\
\bottomrule
\end{tabularx}

\textbf{Response 200:}
\begin{lstlisting}
{
  "accessToken":  "<jwt>",
  "refreshToken": "<jwt>",
  "expireAt":     "30m",
  "basePath":     "john_doe"
}
\end{lstlisting}

\vspace{0.5em}
\subsubsection{\httpget{/v1/auth/me}}

Returns the currently authenticated account.

\textbf{Auth required:} Yes

\textbf{Response 200:}
\begin{lstlisting}
{
  "id":       "<uuid>",
  "username": "john.doe@gmail.com",
  "fullname": "John Doe",
  "email":    "john.doe@gmail.com",
  "provider": "google",
  "rbacId":   "<uuid|null>",
  "rbacName": "admin",
  "rbac":     { ... },
  "homePath": "john_doe"
}
\end{lstlisting}

\vspace{0.5em}
\subsubsection{\httpget{/v1/auth/users}}

Lists all accounts in the system (excluding the caller).

\textbf{Auth required:} Yes

\textbf{Response 200:}
\begin{lstlisting}
[
  {
    "id":       "<uuid>",
    "username": "jane.doe@gmail.com",
    "fullname": "Jane Doe",
    "email":    "jane.doe@gmail.com"
  }
]
\end{lstlisting}

\vspace{0.5em}
\subsubsection{\httpget{/v1/auth/refresh}}

Issues a new \texttt{accessToken} and \texttt{refreshToken}, creating a new session entry.

\textbf{Auth required:} Yes (existing valid session)

\textbf{Response 200:}
\begin{lstlisting}
{
  "accessToken":  "<jwt>",
  "refreshToken": "<jwt>",
  "expireAt":     "30m"
}
\end{lstlisting}

\vspace{0.5em}
\subsubsection{\httpget{/v1/auth/logout}}

Invalidates all active sessions for the authenticated user.

\textbf{Auth required:} Yes

\textbf{Response 200:}
\begin{lstlisting}
{
  "statusCode": 200,
  "messages":   ["Logged out!"]
}
\end{lstlisting}

\newpage

% ══════════════════════════════════════════════════════════════════════════════
\section{Folder Management}
% ══════════════════════════════════════════════════════════════════════════════

Folders form a recursive tree structure in the database. Each folder belongs to an account and optionally has a parent folder. The \texttt{source} JSON field stores the FTP coordinates (host, port, remote path). Folder deletion is hard-delete: the folder and all its descendants are removed from both the database and the FTP server.

\subsection{Endpoints}

\subsubsection{\httppost{/v1/folder}}

Creates a new folder.

\textbf{Auth required:} Yes

\textbf{Request Body (JSON):}

\begin{tabularx}{\linewidth}{llcX}
\toprule
\textbf{Field} & \textbf{Type} & \textbf{Req?} & \textbf{Description} \\
\midrule
\texttt{folderName} & string  & \required & Display name of the folder \\
\texttt{parentId}   & string (UUID) & \optional & Parent folder UUID. Root if omitted \\
\texttt{label}      & string[]& \optional & Up to 20 classification tags \\
\texttt{siteId}     & number  & \required & FTP site identifier (positive integer) \\
\bottomrule
\end{tabularx}

\textbf{Response 201:}
\begin{lstlisting}
{
  "statusCode": 201,
  "messages":   ["Folder Created"],
  "payload":    { "remotePath": "john_doe/projects/my-folder" }
}
\end{lstlisting}

\vspace{0.5em}
\subsubsection{\httpput{/v1/folder/:id}}

Renames or moves an existing folder. Cascades the physical rename on the FTP server and updates all child paths.

\textbf{Auth required:} Yes

\textbf{Path Parameters:}

\begin{tabularx}{\linewidth}{llX}
\toprule
\textbf{Param} & \textbf{Type} & \textbf{Description} \\
\midrule
\texttt{id} & UUID & Folder identifier \\
\bottomrule
\end{tabularx}

\textbf{Request Body (JSON):}

\begin{tabularx}{\linewidth}{llcX}
\toprule
\textbf{Field} & \textbf{Type} & \textbf{Req?} & \textbf{Description} \\
\midrule
\texttt{folderName} & string & \required & New folder name \\
\texttt{parentId}   & string (UUID) & \optional & New parent folder UUID \\
\texttt{label}      & string[] & \optional & Updated labels (max 20) \\
\texttt{siteId}     & number & \required & FTP site identifier \\
\bottomrule
\end{tabularx}

\textbf{Response 201:}
\begin{lstlisting}
{
  "statusCode": 201,
  "messages":   ["Folder Updated"],
  "payload":    { "lastWorkDir": "...", "newWorkDir": "..." }
}
\end{lstlisting}

\vspace{0.5em}
\subsubsection{\httpdelete{/v1/folder/:id}}

Hard-deletes a folder and all its descendants from both the database and the FTP server.

\textbf{Auth required:} Yes

\textbf{Path Parameters:}

\begin{tabularx}{\linewidth}{llX}
\toprule
\textbf{Param} & \textbf{Type} & \textbf{Description} \\
\midrule
\texttt{id} & UUID & Folder identifier \\
\bottomrule
\end{tabularx}

\textbf{Response 200:}
\begin{lstlisting}
{
  "statusCode": 200,
  "messages":   ["Folder Deleted"]
}
\end{lstlisting}

\vspace{0.5em}
\subsubsection{\httpget{/v1/folders}}

Returns a paginated list of all folders.

\textbf{Auth required:} Yes

\textbf{Query Parameters:}

\begin{tabularx}{\linewidth}{llcX}
\toprule
\textbf{Param} & \textbf{Type} & \textbf{Req?} & \textbf{Description} \\
\midrule
\texttt{page}     & number & \optional & Page number (default: 1) \\
\texttt{pageSize} & number & \optional & Items per page (default: 25) \\
\texttt{keyword}  & string & \optional & Case-insensitive name filter \\
\bottomrule
\end{tabularx}

\textbf{Response 200:}
\begin{lstlisting}
{
  "items": [
    {
      "id":         "<uuid>",
      "folderName": "Projects",
      "parentId":   null,
      "accountId":  "<uuid>",
      "label":      ["work", "2026"],
      "createdAt":  "2026-06-17T00:00:00.000Z",
      "updatedAt":  null,
      "source":     { "ftpHost": "...", "ftpPort": 1, "remotePath": "..." }
    }
  ],
  "pagination": { "page": 1, "pageSize": 25, "total": 42 },
  "rbac":       { ... }
}
\end{lstlisting}

\vspace{0.5em}
\subsubsection{\httpget{/v1/folder/:id}}

Retrieves a single folder by its UUID.

\textbf{Auth required:} Yes

\textbf{Response 200:} Same shape as a single item in the \texttt{/v1/folders} list.

\vspace{0.5em}
\subsubsection{\httpget{/v1/folder/:id/real-path}}

Returns the resolved FTP path string for the folder (walks the parent chain).

\textbf{Auth required:} Yes

\textbf{Response 200:}
\begin{lstlisting}
{
  "statusCode": 200,
  "payload":    "parent-name/child-name/leaf-name"
}
\end{lstlisting}

\vspace{0.5em}
\subsubsection{\httpget{/v1/my-folders}}

Returns all folders owned by the authenticated user (no pagination).

\textbf{Auth required:} Yes

\textbf{Response 200:} Array of folder objects (same shape as list items).

\newpage

% ══════════════════════════════════════════════════════════════════════════════
\section{File Management}
% ══════════════════════════════════════════════════════════════════════════════

File records in the database map to physical files stored on an FTP server. Every file belongs to exactly one folder. The unique constraint is \texttt{(folderId, fileName)}. Version history is recorded on every rename/move operation.

\subsection{Endpoints}

\subsubsection{\httpput{/v1/file/:id}}

Renames or moves an existing file. Performs the FTP rename first, then updates the database; rolls back the FTP rename if the DB transaction fails.

\textbf{Auth required:} Yes

\textbf{Path Parameters:}

\begin{tabularx}{\linewidth}{llX}
\toprule
\textbf{Param} & \textbf{Type} & \textbf{Description} \\
\midrule
\texttt{id} & UUID & File identifier \\
\bottomrule
\end{tabularx}

\textbf{Request Body (JSON):}

\begin{tabularx}{\linewidth}{llcX}
\toprule
\textbf{Field} & \textbf{Type} & \textbf{Req?} & \textbf{Description} \\
\midrule
\texttt{folderId}  & string & \required & Target folder UUID \\
\texttt{fileName}  & string & \required & New file name \\
\texttt{siteId}    & number & \required & FTP site identifier \\
\bottomrule
\end{tabularx}

\textbf{Response 201:}
\begin{lstlisting}
{
  "statusCode": 201,
  "messages":   ["File renamed"],
  "payload":    { "lastWorkDir": "...", "newWorkDir": "..." }
}
\end{lstlisting}

\vspace{0.5em}
\subsubsection{\httpdelete{/v1/file/:id}}

Hard-deletes a file from the FTP server and removes all related database records (file, sharing links, version history).

\textbf{Auth required:} Yes

\textbf{Path Parameters:}

\begin{tabularx}{\linewidth}{llX}
\toprule
\textbf{Param} & \textbf{Type} & \textbf{Description} \\
\midrule
\texttt{id} & UUID & File identifier \\
\bottomrule
\end{tabularx}

\textbf{Response 200:}
\begin{lstlisting}
{
  "statusCode": 200,
  "messages":   ["File Deleted"]
}
\end{lstlisting}

\vspace{0.5em}
\subsubsection{\httpget{/v1/files}}

Returns a paginated list of files. Supports filtering by keyword, date range, and owner.

\textbf{Auth required:} Yes

\textbf{Query Parameters:}

\begin{tabularx}{\linewidth}{llcX}
\toprule
\textbf{Param} & \textbf{Type} & \textbf{Req?} & \textbf{Description} \\
\midrule
\texttt{page}      & number & \optional & Page number (default: 1) \\
\texttt{pageSize}  & number & \optional & Items per page (default: 25) \\
\texttt{keyword}   & string & \optional & Case-insensitive file name filter \\
\texttt{startDate} & ISO date & \optional & Created-at lower bound (inclusive) \\
\texttt{endDate}   & ISO date & \optional & Created-at upper bound (exclusive) \\
\texttt{accountId} & UUID & \optional & Filter by owner account \\
\bottomrule
\end{tabularx}

\textbf{Response 200:}
\begin{lstlisting}
{
  "items": [
    {
      "id":        "<uuid>",
      "fileName":  "report.pdf",
      "fileSize":  204800,
      "fileType":  "application/pdf",
      "folderId":  "<uuid>",
      "accountId": "<uuid>",
      "createdAt": "2026-06-17T00:00:00.000Z",
      "updatedAt": null,
      "source":    { "ftpHost": "...", "ftpPort": 1, "remotePath": "..." },
      "folder":    { "folderName": "Documents" },
      "account":   { "username": "...", "fullname": "..." },
      "fileSharings": [ { "id": "...", "account": { ... } } ],
      "action":    { "isUpdate": true, "isDelete": true, "isSharing": true }
    }
  ],
  "pagination": { "page": 1, "pageSize": 25, "total": 10 },
  "rbac":       { ... }
}
\end{lstlisting}

\vspace{0.5em}
\subsubsection{\httpget{/v1/file/:id}}

Retrieves a single file by UUID.

\textbf{Auth required:} Yes

\textbf{Response 200:} Same shape as a single item in the \texttt{/v1/files} list, with sharing targets included.

\vspace{0.5em}
\subsubsection{\httpget{/v1/my-files/:id}}

Lists all files inside the specified folder that belong to the authenticated user. Supports keyword and date range filters via query string (\texttt{keyword}, \texttt{startDate}, \texttt{endDate}).

\textbf{Auth required:} Yes

\textbf{Path Parameters:}

\begin{tabularx}{\linewidth}{llX}
\toprule
\textbf{Param} & \textbf{Type} & \textbf{Description} \\
\midrule
\texttt{id} & UUID & Folder identifier \\
\bottomrule
\end{tabularx}

\textbf{Response 200:} Array of file objects with an added boolean \texttt{isUpdate} flag.

\vspace{0.5em}
\subsubsection{\httpget{/v1/file/history/:id}}

Returns the complete version history of a file.

\textbf{Auth required:} Yes

\textbf{Path Parameters:}

\begin{tabularx}{\linewidth}{llX}
\toprule
\textbf{Param} & \textbf{Type} & \textbf{Description} \\
\midrule
\texttt{id} & UUID & File identifier \\
\bottomrule
\end{tabularx}

\textbf{Response 200:}
\begin{lstlisting}
[
  {
    "id":        "<uuid>",
    "version":   1,
    "json":      { "folderId": "...", "fileName": "old-name.pdf", "siteId": 1 },
    "createdAt": "2026-06-01T00:00:00.000Z",
    "account":   { "username": "...", "fullname": "..." }
  }
]
\end{lstlisting}

\newpage

% ══════════════════════════════════════════════════════════════════════════════
\section{Media Upload \& Streaming}
% ══════════════════════════════════════════════════════════════════════════════

The media module handles the actual binary transfer between the client and the FTP server. It supports single-file upload, bulk folder upload, and streaming download.

\subsection{Endpoints}

\subsubsection{\httppost{/v1/media/upload/folder}}

Bulk-uploads files for an entire folder hierarchy. The system creates all necessary parent directories on the FTP server and then calls \texttt{newFile} for each file. Uses \texttt{multipart/form-data}.

\textbf{Auth required:} Yes

\textbf{Query Parameters:}

\begin{tabularx}{\linewidth}{llcX}
\toprule
\textbf{Param} & \textbf{Type} & \textbf{Req?} & \textbf{Description} \\
\midrule
\texttt{siteId}   & string & \required & FTP site identifier \\
\texttt{folderId} & string (UUID) & \optional & Target folder UUID. Root if omitted \\
\bottomrule
\end{tabularx}

\textbf{Request Body:} \texttt{multipart/form-data} — one or more files in the \texttt{files} field, preserving relative directory paths in file names.

\textbf{Response 200:}
\begin{lstlisting}
{
  "statusCode": 200,
  "messages":   ["Folder uploaded"],
  "payload":    { "uploaded": 5, "failed": 0 }
}
\end{lstlisting}

\vspace{0.5em}
\subsubsection{\httppost{/v1/media/upload/:id}}

Uploads a single file to the folder identified by \texttt{:id}. Saves the binary to the FTP server, then upserts the file record in the database.

\textbf{Auth required:} Yes

\textbf{Path Parameters:}

\begin{tabularx}{\linewidth}{llX}
\toprule
\textbf{Param} & \textbf{Type} & \textbf{Description} \\
\midrule
\texttt{id} & UUID & Destination folder identifier \\
\bottomrule
\end{tabularx}

\textbf{Request Body:} \texttt{multipart/form-data} — single file in the \texttt{file} field.

\textbf{Response 200:}
\begin{lstlisting}
{
  "statusCode": 200,
  "messages":   ["File uploaded"],
  "payload":    { "remotePath": "john_doe/documents" }
}
\end{lstlisting}

\vspace{0.5em}
\subsubsection{\httppost{/v1/media/stream}}

Streams a file from the FTP server to the client. The response is a binary stream with the appropriate \texttt{Content-Type} and \texttt{Content-Disposition} headers.

\textbf{Auth required:} Yes

\textbf{Request Body (JSON):}

\begin{tabularx}{\linewidth}{llcX}
\toprule
\textbf{Field} & \textbf{Type} & \textbf{Req?} & \textbf{Description} \\
\midrule
\texttt{fileName}   & string & \required & Name of the file to stream \\
\texttt{remotePath} & string & \required & FTP directory path \\
\texttt{site}       & number & \required & FTP site identifier \\
\bottomrule
\end{tabularx}

\textbf{Response:} Binary file stream with MIME type resolved from extension.

\vspace{0.5em}
\subsubsection{\httpget{/v1/media/site}}

Returns the list of available FTP site configurations visible to the authenticated user.

\textbf{Auth required:} Yes

\textbf{Response 200:}
\begin{lstlisting}
{
  "statusCode": 200,
  "payload":    [
    { "siteId": 1, "label": "Primary Site" },
    { "siteId": 2, "label": "Backup Site" }
  ]
}
\end{lstlisting}

\newpage

% ══════════════════════════════════════════════════════════════════════════════
\section{Sharing}
% ══════════════════════════════════════════════════════════════════════════════

Resources can be shared between accounts with configurable permissions and optional expiry. Sharing records include a generated link for direct access. Two independent sharing models exist: \emph{file sharing} and \emph{folder sharing}.

\begin{warningbox}
  Sharing permission propagation follows the direction defined in the business rules:
  folder sharing impacts bottom-up (ancestors are also accessible), while file sharing impacts top-down (the parent folder path becomes accessible).
\end{warningbox}

\subsection{File Sharing}

\subsubsection{\httppost{/v1/sharing/file}}

Creates a new file-sharing record.

\textbf{Auth required:} Yes

\textbf{Request Body (JSON):}

\begin{tabularx}{\linewidth}{llcX}
\toprule
\textbf{Field} & \textbf{Type} & \textbf{Req?} & \textbf{Description} \\
\midrule
\texttt{fileId}      & string (UUID) & \required & File to share \\
\texttt{toAccountId} & string (UUID) & \required & Recipient account UUID \\
\texttt{permission}  & enum & \optional & \texttt{READ\_ONLY} (default) or \texttt{READ\_WRITE} \\
\texttt{expiredAt}   & ISO date & \required & Sharing expiry date/time \\
\bottomrule
\end{tabularx}

\textbf{Response 201:}
\begin{lstlisting}
{
  "statusCode":    201,
  "messages":      ["File sharing created"],
  "payload":       {
    "id":            "<uuid>",
    "generatedLink": "https://..."
  }
}
\end{lstlisting}

\vspace{0.5em}
\subsubsection{\httpdelete{/v1/sharing/file/:id}}

Removes a file-sharing record.

\textbf{Auth required:} Yes

\textbf{Response 200:}
\begin{lstlisting}
{ "statusCode": 200, "messages": ["File sharing removed"] }
\end{lstlisting}

\vspace{0.5em}
\subsubsection{\httpget{/v1/sharing/file/:id}}

Retrieves details of a file-sharing record.

\textbf{Auth required:} Yes

\textbf{Response 200:}
\begin{lstlisting}
{
  "id":            "<uuid>",
  "fileId":        "<uuid>",
  "toAccountId":   "<uuid>",
  "permission":    "READ_ONLY",
  "expiredAt":     "2026-12-31T23:59:59.000Z",
  "generatedLink": "https://...",
  "createdAt":     "2026-06-17T00:00:00.000Z"
}
\end{lstlisting}

\subsection{Folder Sharing}

\subsubsection{\httppost{/v1/sharing/folder}}

Creates a new folder-sharing record.

\textbf{Auth required:} Yes

\textbf{Request Body (JSON):}

\begin{tabularx}{\linewidth}{llcX}
\toprule
\textbf{Field} & \textbf{Type} & \textbf{Req?} & \textbf{Description} \\
\midrule
\texttt{folderId}    & string (UUID) & \required & Folder to share \\
\texttt{toAccountId} & string (UUID) & \required & Recipient account UUID \\
\texttt{permission}  & enum & \optional & \texttt{READ\_ONLY} (default) or \texttt{READ\_WRITE} \\
\texttt{expiredAt}   & ISO date & \required & Sharing expiry date/time \\
\bottomrule
\end{tabularx}

\textbf{Response 201:}
\begin{lstlisting}
{
  "statusCode":    201,
  "messages":      ["Folder sharing created"],
  "payload":       {
    "id":            "<uuid>",
    "generatedLink": "https://..."
  }
}
\end{lstlisting}

\vspace{0.5em}
\subsubsection{\httpdelete{/v1/sharing/folder/:id}}

Removes a folder-sharing record.

\textbf{Auth required:} Yes

\textbf{Response 200:}
\begin{lstlisting}
{ "statusCode": 200, "messages": ["Folder sharing removed"] }
\end{lstlisting}

\vspace{0.5em}
\subsubsection{\httpget{/v1/sharing/folder/:id}}

Retrieves details of a folder-sharing record.

\textbf{Auth required:} Yes

\textbf{Response 200:} Same shape as file-sharing detail but with \texttt{folderId} instead of \texttt{fileId}.

\newpage

% ══════════════════════════════════════════════════════════════════════════════
\section{Debug Endpoints}
% ══════════════════════════════════════════════════════════════════════════════

\begin{warningbox}
  Debug endpoints are \textbf{blocked in production}. Requests to these routes when \texttt{NODE\_ENV=production} are silently redirected to \texttt{/}.
\end{warningbox}

\subsection{Endpoints}

\subsubsection{\httppost{/v1/debug/ftp/:siteId}}

Opens an FTP connection to the specified site and checks whether a remote path exists.

\textbf{Path Parameters:}

\begin{tabularx}{\linewidth}{llX}
\toprule
\textbf{Param} & \textbf{Type} & \textbf{Description} \\
\midrule
\texttt{siteId} & integer & Numeric FTP site identifier \\
\bottomrule
\end{tabularx}

\textbf{Request Body (JSON):}

\begin{tabularx}{\linewidth}{llcX}
\toprule
\textbf{Field} & \textbf{Type} & \textbf{Req?} & \textbf{Description} \\
\midrule
\texttt{remotePath} & string & \required & Remote path to inspect on the FTP server \\
\bottomrule
\end{tabularx}

\textbf{Response 200:} Raw FTP directory listing for the given path.

\vspace{0.5em}
\subsubsection{\httpget{/v1/debug/orphans}}

Runs a reconciliation check and returns a summary of orphaned records — database entries with no corresponding FTP artefact, and FTP artefacts with no database record.

\textbf{Response 200:}
\begin{lstlisting}
{
  "summary": {
    "dbFoldersMissingOnFtp":  3,
    "ftpFoldersMissingInDb":  1,
    "dbFilesMissingOnFtp":    2,
    "ftpFilesMissingInDb":    0
  },
  "orphans": { ... }
}
\end{lstlisting}

\vspace{0.5em}
\subsubsection{\httpdelete{/v1/debug/orphans}}

Runs the reconciliation check and \textbf{deletes all orphans}: marks orphaned DB records as \texttt{DEAD} and removes orphaned FTP artefacts.

\textbf{Response 200:}
\begin{lstlisting}
{
  "deleted": {
    "dbFolders":  3,
    "dbFiles":    2,
    "ftpFolders": 1,
    "ftpFiles":   0
  }
}
\end{lstlisting}

\newpage

% ══════════════════════════════════════════════════════════════════════════════
\section{Data Models}
% ══════════════════════════════════════════════════════════════════════════════

All models live in the PostgreSQL \texttt{drive} schema. Timestamps are stored in UTC.

\subsection{Enumerations}

\begin{table}[h!]
\centering
\begin{tabularx}{\linewidth}{llX}
\toprule
\textbf{Enum} & \textbf{Values} & \textbf{Description} \\
\midrule
\texttt{RecordStatus}   & \texttt{ACTIVE}, \texttt{NOT\_ACTIVE}, \texttt{DEAD} & Soft-delete lifecycle for most entities \\
\texttt{SharePermission} & \texttt{READ\_ONLY}, \texttt{READ\_WRITE} & Access level granted in sharing records \\
\bottomrule
\end{tabularx}
\caption{Database enumerations}
\end{table}

\subsection{Account}

Represents an authenticated user. Created or updated on first Google OAuth login.

\begin{longtable}{llD}
\toprule
\textbf{Column} & \textbf{Type} & \textbf{Description} \\
\midrule
\texttt{id}           & UUID (PK)     & Primary key \\
\texttt{dirId}        & UUID?         & Linked directorate \\
\texttt{username}     & string (UNIQUE) & Google email address \\
\texttt{fullname}     & string?       & Display name from Google profile \\
\texttt{email}        & string?       & Email (may differ from username) \\
\texttt{provider}     & string        & OAuth provider (\texttt{google}) \\
\texttt{rbacId}       & UUID?         & Assigned RBAC role \\
\texttt{accountInfo}  & JSON?         & Raw Google profile data \\
\texttt{quota}        & int           & Storage quota in bytes (default 0) \\
\texttt{createdAt}    & DateTime      & Account creation timestamp \\
\texttt{updatedAt}    & DateTime?     & Last update timestamp \\
\texttt{recordStatus} & RecordStatus  & Lifecycle status (default \texttt{NOT\_ACTIVE}) \\
\bottomrule
\caption{Account model}
\end{longtable}

\subsection{Session}

Stores active JWT tokens. Each login or refresh creates a new session record.

\begin{longtable}{llD}
\toprule
\textbf{Column} & \textbf{Type} & \textbf{Description} \\
\midrule
\texttt{id}           & UUID (PK)    & Primary key \\
\texttt{accountId}    & UUID (FK)    & Owner account \\
\texttt{appCode}      & string       & Application identifier (default \texttt{main\_app}) \\
\texttt{jwtHash}      & string       & The raw JWT access token (used for replay prevention) \\
\texttt{recordStatus} & RecordStatus & \texttt{ACTIVE} = valid session \\
\texttt{createdAt}    & DateTime     & Creation timestamp \\
\bottomrule
\caption{Session model}
\end{longtable}

\subsection{Folder}

Represents a directory in the virtual drive, backed by a physical FTP directory.

\begin{longtable}{llD}
\toprule
\textbf{Column} & \textbf{Type} & \textbf{Description} \\
\midrule
\texttt{id}           & UUID (PK)    & Primary key \\
\texttt{parentId}     & UUID?        & Self-referential parent (null = root) \\
\texttt{accountId}    & UUID (FK)    & Owner account \\
\texttt{folderName}   & string       & Display/directory name \\
\texttt{source}       & JSON         & FTP coordinates: \texttt{ftpHost}, \texttt{ftpPort}, \texttt{remotePath} \\
\texttt{label}        & JSON[]       & Classification tags \\
\texttt{createdAt}    & DateTime     & Creation timestamp \\
\texttt{updatedAt}    & DateTime?    & Last update timestamp \\
\texttt{recordStatus} & RecordStatus & Lifecycle status \\
\bottomrule
\caption{Folder model}
\end{longtable}

\subsection{File}

Represents a single file. Unique per \texttt{(folderId, fileName)}.

\begin{longtable}{llD}
\toprule
\textbf{Column} & \textbf{Type} & \textbf{Description} \\
\midrule
\texttt{id}           & UUID (PK)    & Primary key \\
\texttt{folderId}     & UUID (FK)    & Containing folder \\
\texttt{accountId}    & UUID (FK)    & Owner account \\
\texttt{fileName}     & string       & File name (unique per folder) \\
\texttt{fileHash}     & string?      & MD5 hash from FTP \texttt{XMD5} command \\
\texttt{fileSize}     & int          & Size in bytes \\
\texttt{fileType}     & string       & MIME type \\
\texttt{source}       & JSON         & FTP coordinates \\
\texttt{createdAt}    & DateTime     & Upload timestamp \\
\texttt{updatedAt}    & DateTime?    & Last modification timestamp \\
\texttt{recordStatus} & RecordStatus & Lifecycle status \\
\bottomrule
\caption{File model}
\end{longtable}

\subsection{FileHistory}

Records every rename/move operation on a file for audit and rollback purposes.

\begin{longtable}{llD}
\toprule
\textbf{Column} & \textbf{Type} & \textbf{Description} \\
\midrule
\texttt{id}        & UUID (PK) & Primary key \\
\texttt{fileId}    & UUID (FK) & Associated file \\
\texttt{accountId} & UUID (FK) & Account that made the change \\
\texttt{json}      & JSON      & Snapshot of the change DTO \\
\texttt{version}   & int       & Monotonically increasing version counter \\
\texttt{createdAt} & DateTime  & Change timestamp \\
\bottomrule
\caption{FileHistory model}
\end{longtable}

\subsection{FolderSharing / FileSharing}

Both sharing models have identical structure, differing only in the resource they reference.

\begin{longtable}{llD}
\toprule
\textbf{Column} & \textbf{Type} & \textbf{Description} \\
\midrule
\texttt{id}            & UUID (PK)       & Primary key \\
\texttt{folderId / fileId} & UUID (FK)   & Shared resource \\
\texttt{accountId}     & UUID (FK)       & Sharing owner \\
\texttt{toAccountId}   & UUID (FK)       & Sharing recipient \\
\texttt{permission}    & SharePermission & \texttt{READ\_ONLY} or \texttt{READ\_WRITE} \\
\texttt{expiredAt}     & DateTime?       & Optional expiry \\
\texttt{generatedLink} & string          & Public access link \\
\texttt{createdAt}     & DateTime        & Creation timestamp \\
\texttt{recordStatus}  & RecordStatus    & Lifecycle status \\
\bottomrule
\caption{Sharing models}
\end{longtable}

\subsection{Trace / TraceSpan}

Observability tables populated by the OpenTelemetry middleware.

\begin{table}[h!]
\centering
\begin{tabularx}{\linewidth}{llX}
\toprule
\textbf{Column} & \textbf{Type} & \textbf{Description} \\
\midrule
\texttt{Trace.id}             & string      & UUID v7 trace ID \\
\texttt{Trace.url}            & string      & Request URL \\
\texttt{Trace.method}         & string      & HTTP method \\
\texttt{Trace.status}         & int         & HTTP response status code \\
\texttt{Trace.responseTimeMs} & int         & Response time in milliseconds \\
\texttt{TraceSpan.traceId}    & string?     & Parent trace ID \\
\texttt{TraceSpan.context}    & string      & Span name/context label \\
\texttt{TraceSpan.durationMs} & int         & Span duration in milliseconds \\
\bottomrule
\end{tabularx}
\caption{Observability models}
\end{table}

\newpage

% ══════════════════════════════════════════════════════════════════════════════
\section{Error Handling}
% ══════════════════════════════════════════════════════════════════════════════

All errors follow a consistent JSON envelope:

\begin{lstlisting}
{
  "errCode":    "FILE_NOT_FOUND",
  "statusCode": 404,
  "messages":   ["Selected file not found!"]
}
\end{lstlisting}

\begin{table}[h!]
\centering
\begin{tabularx}{\linewidth}{llX}
\toprule
\textbf{errCode} & \textbf{HTTP} & \textbf{Meaning} \\
\midrule
\texttt{UNAUTHORIZED}       & 401 & Missing, invalid, or expired JWT/session \\
\texttt{GUARD\_ERROR}       & 500 & Unexpected error in authentication middleware \\
\texttt{VALIDATION\_ERROR}  & 422 & Request DTO validation failed \\
\texttt{ACCOUNT\_NOT\_FOUND}& 404 & Account referenced does not exist \\
\texttt{FOLDER\_NOT\_FOUND} & 404 & Folder not found or not owned by caller \\
\texttt{FILE\_NOT\_FOUND}   & 404 & File not found or not owned by caller \\
\texttt{FILE\_EXIST}        & 409 & Duplicate file name in the same folder \\
\texttt{SOURCE\_NOT\_FOUND} & 404 & File record has no FTP source metadata \\
\bottomrule
\end{tabularx}
\caption{Known error codes}
\end{table}

\newpage

% ══════════════════════════════════════════════════════════════════════════════
\section{Environment Configuration}
% ══════════════════════════════════════════════════════════════════════════════

The application reads configuration from environment variables validated by \texttt{envalid}. All variables marked \textbf{required} have no default and must be set before startup.

\begin{longtable}{llcE}
\toprule
\textbf{Variable} & \textbf{Type} & \textbf{Req?} & \textbf{Description} \\
\midrule
\texttt{APP\_NAME}                & string  & \optional & Application name (default: \texttt{DriveAPI}) \\
\texttt{NODE\_ENV}                & enum    & \optional & \texttt{development} | \texttt{production} | \texttt{staging} | \texttt{test} \\
\texttt{NODE\_REJECT\_UNAUTHORIZE}& bool    & \optional & Reject unauthorised TLS certs (default: \texttt{false}) \\
\texttt{LOG}                      & enum    & \optional & \texttt{all} | \texttt{none} | \texttt{prisma} | \texttt{ftp} (default: \texttt{none}) \\
\texttt{PORT}                     & number  & \optional & HTTP listen port (default: \texttt{9000}) \\
\texttt{API\_URL}                 & string  & \optional & Public base URL (default: \texttt{localhost}) \\
\texttt{JWT\_SECRET}              & string  & \optional & HS256 signing secret (default: \texttt{jwt-secret} — \textbf{change in prod!}) \\
\texttt{JWT\_EXPIRE}              & string  & \optional & Token lifetime (default: \texttt{30m}) \\
\texttt{E\_SALT}                  & string  & \required & Encryption salt for FTP credentials \\
\texttt{DATABASE\_URL}            & string  & \optional & PostgreSQL connection string \\
\texttt{FTP\_HOST}                & string  & \optional & FTP server hostname (default: \texttt{localhost}) \\
\texttt{FTP\_USERNAME}            & string  & \required & FTP authentication username \\
\texttt{FTP\_PASSWORD}            & string  & \required & FTP authentication password \\
\texttt{FTP\_CONFIG}              & mixed   & \optional & FTP connection mode: \texttt{implicit} (default) | \texttt{explicit} | \texttt{true}/\texttt{false} \\
\texttt{GOOGLE\_CLIENT\_ID}       & string  & \required & Google OAuth 2.0 client ID \\
\texttt{GOOGLE\_CLIENT\_SECRET}   & string  & \required & Google OAuth 2.0 client secret \\
\texttt{GOOGLE\_REDIRECT\_URL}    & string  & \required & Registered OAuth callback URL \\
\bottomrule
\caption{Environment variables}
\end{longtable}

\newpage

% ══════════════════════════════════════════════════════════════════════════════
\section{Background Jobs}
% ══════════════════════════════════════════════════════════════════════════════

Background tasks are handled by \textbf{Inngest}, an event-driven job platform. The Inngest serve route is registered at \texttt{/v1/api/inngest}.

\subsection{Orphan Reconciliation}

The \texttt{check-orphans} job compares every folder and file record in the database against the actual FTP directory tree:

\begin{itemize}
  \item \textbf{DB folders missing on FTP} — records that have no matching directory on any FTP site.
  \item \textbf{FTP folders missing in DB} — directories on the FTP server with no corresponding database record.
  \item \textbf{DB files missing on FTP} — file records whose physical FTP file has been deleted.
  \item \textbf{FTP files missing in DB} — FTP files with no database record.
\end{itemize}

The \texttt{reconcile-orphaned-files} Inngest function can be triggered to automatically clean up these discrepancies.

\newpage

% ══════════════════════════════════════════════════════════════════════════════
\section{Middleware Pipeline}
% ══════════════════════════════════════════════════════════════════════════════

Requests pass through the following middleware layers in order:

\begin{enumerate}[label=\arabic*.]
  \item \textbf{Context Storage} (\texttt{hono/context-storage}) — makes the Hono context available globally within a request.
  \item \textbf{Trace ID injection} — generates a UUID v7 trace ID and sets the \texttt{x-trace-id} response header.
  \item \textbf{CORS} (\texttt{hono/cors}) — enables cross-origin resource sharing.
  \item \textbf{Telemetry} (\texttt{useTelemetry}) — records request/response metrics and spans to the \texttt{Trace} and \texttt{TraceSpan} tables via OpenTelemetry.
  \item \textbf{Secure Headers} (\texttt{hono/secure-headers}) — adds OWASP-recommended security headers.
  \item \textbf{Pretty JSON} (\texttt{hono/pretty-json}) — formats JSON responses when the client sends \texttt{?pretty=1}.
  \item \textbf{Request Validator} (\texttt{Validate.for}) — validates the request body, query, or path params against the specified DTO class using \texttt{class-validator}.
  \item \textbf{Auth Guard} (\texttt{Guard.validate}) — verifies the JWT and active session; injects the \texttt{account} context variable.
  \item \textbf{RBAC Validator} (\texttt{rbac.validator}) — optional role-based access control check.
  \item \textbf{Error Handler} (\texttt{errorHandler}) — catches \texttt{HttpException} instances and all uncaught errors, returning a structured JSON error response.
\end{enumerate}

\newpage

% ══════════════════════════════════════════════════════════════════════════════
\section{Endpoint Summary}
% ══════════════════════════════════════════════════════════════════════════════

\begin{longtable}{lllF}
\toprule
\textbf{Method} & \textbf{Path} & \textbf{Auth} & \textbf{Description} \\
\midrule
\multicolumn{4}{l}{\textit{OAuth \& Auth}} \\
\midrule
GET    & /v1/oauth/google              & No  & Initiate Google OAuth flow \\
GET    & /v1/oauth/google/callback     & No  & Google OAuth callback \\
GET    & /v1/auth/me                   & Yes & Get current user \\
GET    & /v1/auth/users                & Yes & List all users \\
GET    & /v1/auth/refresh              & Yes & Refresh access token \\
GET    & /v1/auth/logout               & Yes & Invalidate session \\
\midrule
\multicolumn{4}{l}{\textit{Folders}} \\
\midrule
POST   & /v1/folder                    & Yes & Create folder \\
PUT    & /v1/folder/:id                & Yes & Rename / move folder \\
DELETE & /v1/folder/:id                & Yes & Hard-delete folder \\
GET    & /v1/folders                   & Yes & List folders (paginated) \\
GET    & /v1/folder/:id                & Yes & Get folder by ID \\
GET    & /v1/folder/:id/real-path      & Yes & Get folder FTP path \\
GET    & /v1/my-folders                & Yes & Get my folders \\
\midrule
\multicolumn{4}{l}{\textit{Files}} \\
\midrule
PUT    & /v1/file/:id                  & Yes & Rename / move file \\
DELETE & /v1/file/:id                  & Yes & Hard-delete file \\
GET    & /v1/files                     & Yes & List files (paginated) \\
GET    & /v1/file/:id                  & Yes & Get file by ID \\
GET    & /v1/my-files/:id              & Yes & List files in my folder \\
GET    & /v1/file/history/:id          & Yes & Get file version history \\
\midrule
\multicolumn{4}{l}{\textit{Media}} \\
\midrule
POST   & /v1/media/upload/folder       & Yes & Bulk folder upload \\
POST   & /v1/media/upload/:id          & Yes & Single file upload \\
POST   & /v1/media/stream              & Yes & Stream / download file \\
GET    & /v1/media/site                & Yes & List FTP sites \\
\midrule
\multicolumn{4}{l}{\textit{Sharing}} \\
\midrule
POST   & /v1/sharing/file              & Yes & Create file sharing \\
DELETE & /v1/sharing/file/:id          & Yes & Remove file sharing \\
GET    & /v1/sharing/file/:id          & Yes & Get file sharing \\
POST   & /v1/sharing/folder            & Yes & Create folder sharing \\
DELETE & /v1/sharing/folder/:id        & Yes & Remove folder sharing \\
GET    & /v1/sharing/folder/:id        & Yes & Get folder sharing \\
\midrule
\multicolumn{4}{l}{\textit{Debug (non-production only)}} \\
\midrule
POST   & /v1/debug/ftp/:siteId         & No  & Test FTP path on site \\
GET    & /v1/debug/orphans             & No  & Inspect orphaned records \\
DELETE & /v1/debug/orphans             & No  & Delete all orphaned records \\
\midrule
\multicolumn{4}{l}{\textit{Docs \& Jobs}} \\
\midrule
GET    & /docs                         & No  & Swagger UI \\
GET    & /docs/openapi.json            & No  & OpenAPI 3.x specification \\
POST   & /v1/api/inngest               & No  & Inngest function serve endpoint \\
\bottomrule
\caption{Complete endpoint summary}
\end{longtable}

\newpage

% ══════════════════════════════════════════════════════════════════════════════
\section{Security Considerations}
% ══════════════════════════════════════════════════════════════════════════════

\begin{enumerate}[label=\arabic*.]
  \item \textbf{JWT Secret} — the default value \texttt{jwt-secret} must be overridden with a cryptographically strong secret in staging and production environments.
  \item \textbf{Session replay prevention} — the raw JWT is stored as \texttt{jwtHash} in the \texttt{Session} table. The guard middleware validates both the JWT signature \emph{and} the existence of an active session entry, preventing use of stolen tokens after logout.
  \item \textbf{Ownership enforcement} — all mutating operations filter by \texttt{accountId = currentUser.id}, ensuring users can only modify their own resources.
  \item \textbf{FTP TLS} — \texttt{NODE\_REJECT\_UNAUTHORIZE=false} by default eases local development; set to \texttt{true} in production to validate FTP server certificates.
  \item \textbf{Encryption salt} — the \texttt{E\_SALT} environment variable must be set; it is used for encrypting FTP credential material at rest.
  \item \textbf{Debug endpoint isolation} — debug routes check \texttt{NODE\_ENV} at request time and redirect to \texttt{/} in production, providing a safety net against accidental exposure.
  \item \textbf{Secure headers} — the \texttt{hono/secure-headers} middleware adds \texttt{X-Content-Type-Options}, \texttt{X-Frame-Options}, \texttt{Strict-Transport-Security}, and other OWASP-recommended headers.
  \item \textbf{Input validation} — every DTO is validated with \texttt{class-validator} before the handler executes; malformed requests are rejected with \texttt{422 Unprocessable Entity}.
\end{enumerate}

% ══════════════════════════════════════════════════════════════════════════════
\section*{Appendix A — Swagger / OpenAPI}
\addcontentsline{toc}{section}{Appendix A — Swagger / OpenAPI}
% ══════════════════════════════════════════════════════════════════════════════

The API ships with an auto-generated OpenAPI 3.x specification. When the service is running, it is available at:

\begin{itemize}
  \item \textbf{Swagger UI:} \texttt{GET /docs}
  \item \textbf{Raw JSON spec:} \texttt{GET /docs/openapi.json}
\end{itemize}

The specification is generated from the swagger definition in \texttt{src/lib/swagger.ts} using \texttt{@hono/swagger-ui}.

% ══════════════════════════════════════════════════════════════════════════════
\section*{Appendix B — Running Locally}
\addcontentsline{toc}{section}{Appendix B — Running Locally}
% ══════════════════════════════════════════════════════════════════════════════

\begin{lstlisting}
# Install dependencies
npm install

# Development (hot-reload via tsx)
npm run dev

# Production build (generates Prisma client + runs migrations + bundles)
npm run build

# Start production server
npm start
\end{lstlisting}

The server listens on \texttt{PORT} (default \texttt{9000}).

\vspace{1em}
\textit{This document was generated on \today.}

\end{document}
